\documentclass[UTF8,a4paper,12pt]{ctexart}

%=====================
% 基础宏包
%=====================
\usepackage{geometry}
\usepackage{fancyhdr}
\usepackage{amsmath,amssymb}
\usepackage{graphicx}
\usepackage{booktabs}
\usepackage{float}
\usepackage{algorithm}
\usepackage{algorithmic}
\usepackage{listings}
\usepackage{xcolor}
\usepackage{hyperref}
\usepackage{setspace}

%=====================
% 页面设置
%=====================
\geometry{
    a4paper,
    left=2.5cm,
    right=2.5cm,
    top=2.5cm, 
    bottom=2.5cm
}

%=====================
% 行距
%=====================
\linespread{1.25}


%=====================
% 页码设置
% 摘要开始编号
%=====================
\pagestyle{fancy}
\fancyhf{}
\fancyfoot[C]{\arabic{page}}

\renewcommand{\headrulewidth}{0pt}


%=====================
% 图片路径
%=====================
\graphicspath{{figures/}}


%=====================
% 代码格式
%=====================
\lstset{
    basicstyle=\ttfamily\small,
    numbers=left,
    numberstyle=\tiny,
    frame=single,
    breaklines=true
}


%=====================
% 标题格式
%=====================
\title{
    \heiti
    全国大学生数学建模竞赛论文
}

\author{}
\date{}


\begin{document}


%%%%%%%%%%%%%%%%%%%%%%%%%%%%%%%%
% 摘要页
% 电子版第一页
%%%%%%%%%%%%%%%%%%%%%%%%%%%%%%%%

\pagenumbering{arabic}

\begin{center}

{\zihao{3}\bfseries 论文题目：XXX模型研究}

\vspace{1cm}

\end{center}


\section*{摘要}

这里填写摘要。

摘要应包含：

\begin{itemize}
    \item 问题分析；
    \item 建模方法；
    \item 求解过程；
    \item 主要结果；
    \item 模型评价。
\end{itemize}


\vspace{0.5cm}


\noindent
\textbf{关键词：}
关键词1；关键词2；关键词3


\newpage


%%%%%%%%%%%%%%%%%%%%%%%%%%%%%%%%
% 正文
%%%%%%%%%%%%%%%%%%%%%%%%%%%%%%%%


\section{问题重述}

描述赛题背景和需要解决的问题。



\section{问题分析}


分析问题特点，说明建模思路。


\section{模型假设}

\begin{enumerate}
    \item 假设1；
    \item 假设2；
    \item 假设3 。
\end{enumerate}



\section{符号说明}


\begin{table}[H]
\centering

\begin{tabular}{ccc}
\toprule
符号 & 含义 & 单位\\
\midrule
$x$ & 变量 & -\\
$y$ & 输出量 & -\\
\bottomrule

\end{tabular}

\caption{符号说明}

\end{table}



\section{模型建立}


数学模型：

\begin{equation}
f(x)=ax^2+bx+c
\end{equation}


\subsection{子模型1}


描述模型建立过程。


\subsection{子模型2}


描述模型建立过程。



\section{模型求解}


算法：

\begin{algorithm}[H]
\caption{求解算法}
\begin{algorithmic}

\STATE 输入数据
\STATE 初始化参数
\STATE 迭代计算
\STATE 输出结果

\end{algorithmic}
\end{algorithm}



\section{结果分析}


\begin{figure}[H]
\centering



\caption{结果示意图}

\end{figure}



\section{模型评价}


\subsection{模型优点}

说明模型优势。$\int_{a}^{b}x$



\subsection{模型缺点}

说明模型不足。


\subsection{改进方向}

说明未来优化。


\section{参考文献}


\bibliographystyle{plain}

\bibliography{references}



%%%%%%%%%%%%%%%%%%%%%%%%%%%%%%%%
% 附录
%%%%%%%%%%%%%%%%%%%%%%%%%%%%%%%%

\appendix

\section{支撑材料文件列表}


例如：

\begin{itemize}
    \item data.xlsx
    \item model.py
    \item result.csv
\end{itemize}



\section{源程序代码}


Python 示例：

\begin{lstlisting}[language=Python]

import numpy as np

def model(x):
    return x*x

print(model(2))

\end{lstlisting}



\end{document}